\begin{problem}[Sheaf of solutions of a differential equation]
Let $M$ be a smooth manifold and $L:C^\infty_M\to\mathcal E$ a local linear differential operator into a sheaf $\mathcal E$. Define $\mathcal S(U)=\{u\in C^\infty(U):Lu=0\}$. Prove that $\mathcal S$ is a sheaf. Compute it for $L=d/dx$ on an interval and for $L=d^2/dx^2$ on an interval.
\end{problem}

\paragraph{Why this problem matters.}
Many important sheaves arise as local solution spaces rather than as all functions.

\paragraph{Useful ideas before starting.}
The kernel condition is local. For the ODE examples, solve the differential equations on each connected interval.

\paragraph{Guided hints.}
For $u=0$, local solutions are affine functions $ax+b$.

\ifdefined\FullProblemDossiers
\begin{solution}
Compatible local solutions glue to a smooth function $u$ because $C^\infty_M$ is a sheaf. On each cover member, $Lu=Lu_i=0$, and since $L$ is local this implies $Lu=0$ globally. Locality is inherited from smooth functions, so $\mathcal S$ is a sheaf. On a connected interval, $u'=0$ gives constant functions, hence the constant sheaf $\underline\RR$. For $u''=0$, sections over a connected interval are affine functions $ax+b$, a two-dimensional vector space; on disconnected opens one chooses such a pair independently on each component.
\end{solution}

\paragraph{What happened mathematically.}
Kernels of local operators naturally satisfy the sheaf axiom, anticipating kernel sheaves in VI/15.

\paragraph{Extensions.}
Study local harmonic functions or locally exact differential forms.

\paragraph{Provenance.}
New canonical synthesis from the standard function-sheaf corpus.
\else
\paragraph{Provenance.}
New canonical synthesis from the standard function-sheaf corpus.
\fi
